\documentclass[10pt,a4paper]{amsart}
\usepackage[margin=19mm]{geometry}
\usepackage{amsmath,amssymb,mathtools}
\usepackage[scr=rsfs,cal=euler]{mathalfa}
\usepackage{xfrac}
\usepackage[unicode,hidelinks,bookmarks=false]{hyperref}
\newcommand{\ie}{i.e.\ }
\newcommand{\cf}{cf.\ }
\newcommand{\resp}{respectively\ }
\newcommand{\Subsection}{\subsection}
\providecommand{\nobreakdash}{\nobreak\hbox{-}}
\setlength{\emergencystretch}{2em}
\multlinegap0pt
\usepackage{amssymb}
\usepackage{enumerate}
\newtheorem{proposition}{Proposition}
\newcommand{\aconn}{$\alpha$-connected}
\newcommand{\aend}{$\alpha$-end}
\newcommand{\akconn}{$(\alpha,K)$-connected}
\newcommand{\akpconn}{$(\alpha,K')$-connected}
\newcommand{\caleinf}{{\mathcal E}^\alpha_\infty}
\newcommand{\calek}{{\mathcal E}^\alpha_K}
\newcommand{\calekp}{{\mathcal E^\alpha}_{K'}}
\newcommand{\ph}{\varphi}
\title{Sparse metric spaces and sparse ends}
\author{William Geller, Michal Misiurewicz}
\date{Source: arXiv:2606.12202v1 (CC BY 4.0)}
\begin{document}
\maketitle

\noindent\emph{Source and licence.} Sparse metric spaces and sparse ends. Version arXiv:2606.12202v1 (CC BY 4.0), distributed by arXiv under the Creative Commons Attribution 4.0 International licence, http://creativecommons.org/licenses/by/4.0/, as recorded in the arXiv licence metadata for this exact version and shown on the abstract page https://arxiv.org/abs/2606.12202v1. Full licence notice: \texttt{licenses/arxiv-2606.12202-CC-BY-4.0.txt}.\par\smallskip

\noindent\textit{Editorial scope.} Counted unit: the article's proposition aendinvariance (source0.tex lines 934--942). The article's complete original proof of it is delivered with it (source0.tex lines 944--983, one \texttt{proof} environment of 40 lines). No mathematics is written, repaired or re-derived: the statement and the proof are the article's own lines, in the article's own notation, and no retained line is substituted or re-flowed.  Retained uncounted as the source-local context the proof needs: lines 775--782.  Nothing was omitted from any retained span, so the package carries no omission record.  No retained line contains a citation, so the wrapper carries no bibliography and no bibliography entry.  No retained line contains a figure environment, an image-inclusion command or drawing code, so no image is delivered and the package declares no asset dependency.  Editorial formatting only, in addition: an AMS \texttt{amsart} preamble that restates the article's own declarations and macros used by the retained lines, namely the environments \texttt{proposition} on separate counters, together with the standard packages those lines need.  The retained environments print in this document as Proposition 1 in order of appearance, whereas the article prints the same items with the numbering of its own full document.  The complete native source of the article as distributed in the arXiv e-print for this exact version is bundled unchanged as \texttt{sources/202609/source0.tex}.
\begin{proposition}\label{aconninvariance}
Quasi-isometries preserve \aconn ness. That is, if $S$ is an \aconn
\ subset of a pointed metric space $(X,x_0,d)$ and $\ph: (X,x_0,d) \to
(X',x'_0,d')$ is a quasi-isometry to another pointed metric space,
then $S'=\ph(S)$ is \aconn. More specifically, if $S$ is \akconn, then
$S'$ is \akpconn, where $K'=\max\{C(K+1), 2^\alpha C^{1+\alpha} K\}$
for a constant $C$ associated with the quasi-isometry.
\end{proposition}

\begin{proposition}\label{aendinvariance}
A quasi-isometry induces a homeomorphism of the spaces of \aend s.

That is, if $\alpha \geq 0$ and $S$ is a subset of a pointed metric
space $(X,x_0,d)$ with \aend s $\caleinf(S)$ and $\ph: (X,x_0,d) \to
(X',x'_0,d')$ is a quasi-isometry to another pointed metric space,
then $\caleinf(S)$ is homeomorphic to $\caleinf(S')$, where
$S'=\ph(S)$.
\end{proposition}

\begin{proof}
We first show that $\ph$ induces a bijection of the set of \aend s. By
symmetry, it suffices to show that $\ph$ induces a map from \aend s of
$S$ to \aend s of $S'$ and then check that this map is one to one.

If $e$ is an \aend \ of $S$, then $e=[E]$, where $E \in \calek(S)$ for
some $K$. Then it follows from Proposition \ref{aconninvariance} that
$\ph$ takes $E$ to some $E' \in \calekp(S')$, where $K'$ is as given
in the proposition. So the induced map $\ph_*$ on \aend s takes $e \in
\caleinf(S)$ to $e' \in \caleinf(S')$, where $e' = [E']$.

Now suppose that $e_1$ and $e_2$ are \aend s of $S$ and
$\ph_*(e_1)=\ph_*(e_2)=e'$. Let $\psi: (X',x'_0,d') \to (X,x_0,d)$ be
an inverse quasi-isometry of $\ph$ with $\psi(S') \subset S$ and
$\psi(S')$ $K_0$-dense in $S$, and let $e= \psi_*(e')$. If $e_1 =
[E_{K_1}]$ and $e_2 = [E_{K_2}]$, and we take $E_K$ with $e = [E_K]$
and $K \geq \max(K_1,K_2)$, then $E_{K_1}$ and $E_{K_2}$ are contained
in $E_K$ so we have that $e_1 = e_2$. So $\ph_*$ is a bijection on the
set of \aend s.

It remains to prove that $\ph_*$ is a homeomorphism. By symmetry, it
suffices to show that $\ph_*$ is continuous. We show that $\ph_*$ is
continuous at every \aend \ $e$ of $S$.

Let $e=[E]$, where $E \in \calek(S)$ for some $K$. Let $e'=\ph_*(e)
=[E']$, where $E' \in \calekp(S')$ for $K' \geq K'_0$, where $K'_0$ is
the constant given in Proposition \ref{aconninvariance}.

A basic neighborhood $V$ of the \aend\ $e'$ corresponds to the
unbounded \akpconn\ component of $S' \setminus B(x'_0, R')$ containing
$e'$ for some $R'$. We find a basic neighborhood $U$ of $e$ mapped by
$\ph$ into $V$. Let $R=CR'+C$, where $C$ is the constant in
Proposition \ref{aconninvariance}. Then $\ph(X \setminus B(x_0, R))
\subset X' \setminus B(x'_0, R')$ since $\ph(B(x_0, R)) \supset
B(x'_0, R')$. So $\ph(S \setminus B(x_0, R)) \subset S' \setminus
B(x'_0, R')$. Let $U$ be the basic neighborhood of $e$ corresponding
to the unbounded \akconn\ component of $S \setminus B(x_0, R)$
containing $e$. Then $\ph(U)$ is an \akpconn\ subset of $S' \setminus
B(x'_0, R')$, and it contains $e'$, so it is contained in $V$.
\end{proof}

\end{document}
